% -----------------------------------------------------------------------------------------------------
% Hypothesis B1: statements restated from theory/B1_preconditioned_saga.tex (Appendix B1). Numbers quoted, on the
% instances of Figure B1 (n = 1000, d = 50, kappa = 1e3, alpha = 1, 2, data seed 0; B1 note, comment after Corollary
% cor:B1-kappa-traj), with R_P^2 the upper bound of eq:B1-logistic-RP as in the experiments:
%   results/B1_summary.json: min_i w_i = sigma_prime_min = 7.39e-3, 6.10e-3; lmin(M_P) = mu_P_local of 'stat' =
%     0.1166, 0.1086 (15.8 and 17.8 times min_i w_i); R_P^2/lmin(M_P) = 166.4, 102.9; R^2/lmin(F''(eta*)) =
%     R2_P/mu_P_local of 'I' = 575.8, 980.9;
%   first bound of eq:B1-kappa-traj-stat with m = min_i w_i: 1735, 1221; upper bound of eq:B1-kappa-traj-I for P = I
%     with the same m: 968.5, 999.3 (recomputed from the data);
%   kappa_eff^P(theta_{t-1}) = R_P^2/mu_eff^P over passes 15--20, geometric mean over the instrumented states and the
%     5 seeds (results/B1.pkl): 84.4, 31.0 (Sigma_hat + lambda I) and 460.5, 877.6 (P = I).
% -----------------------------------------------------------------------------------------------------
\subsection{Preconditioned SAGA}\label{sec:th-B1}

Hypothesis B1 runs SAGA in a fixed quadratic metric to see which steps of the proof of Proposition~\ref{prop:59} are
too loose, and asks whether the metric improves the algorithm or only its analysis. Let $P$ be a fixed symmetric
positive definite matrix, $\operatorname{cond}(P)=\lmax(P)/\lmin(P)$ its condition number, $\norm{v}_P=(v^\top Pv)^{1/2}$
the norm of parameters and $\norm{g}_{P^{-1}}=(g^\top P^{-1}g)^{1/2}$ the dual norm of gradients and memories. The
algorithm of \resref{def:B1-PSAGA} keeps the direction $\omega_t$ and the memory of SAGA~\eqref{eq:saga} and, by
\resref{lem:B1-step}, steps $\theta_t=\theta_{t-1}-\gamma P^{-1}\omega_t$; it is Bregman-SAGA~\cite{dragomir2021fast}
with the reference function $h(\theta)=\frac12\theta^\top P\theta$. A linear change of variables reduces it to SAGA:
\restate{prop:B1-transfer}
Every step of the proof of Proposition~\ref{prop:59} thus holds for $P$-SAGA, in the norms $\norm{\cdot}_P$ and
$\norm{\cdot}_{P^{-1}}$ and with the constants of the transformed functions, written here for $C^2$ functions (for
$C^1$ functions, see \resref*{rem:B1-C1}):
\restate{def:B1-kappaP}
Proposition~\ref{prop:59} then transfers verbatim, with these constants in place of $(R^2,\mu)$:
\restate{cor:B1-rate}
The metric acts on the rate only through $\kappaP$; $\operatorname{cond}(P)$ enters only the constant of the Euclidean
bound. By \resref{lem:B1-compare}, $\kappa/\operatorname{cond}(P)\le\kappaP\le\kappa\operatorname{cond}(P)$, so that a
fixed metric changes the certified rate by at most the factor $\operatorname{cond}(P)$, in either direction. In the
proof written in the $P$-metric (Section~\ref{sec:B1-proofP}), I1 and I3 change only their norm, I2 and I4 their norm
and their constant ($R_P^2$, $\mu_P$), I5 depends on $P$ only through $3\mu_P/(16R_P^2)$, and I6 through both norms and
$R_P^4$. The tightness ratios of I2 and I4 in the $P$-geometry (\resref*{def:B1-ratios}) are
$r^P_{\rm coco}(t)=\LEff^P(\theta_{t-1})/R_P^2$ and $r^P_{\rm sc}(t)=\mu_P/\muEff^P(\theta_{t-1})$, with the effective
constants $\muEff^P(\theta)=\inner{\Fp(\theta)}{\theta-\etastar}/\norm{\theta-\etastar}_P^2$ and
$\LEff^P(\theta)=\frac1n\sum_i\norm{\fip(\theta)-\fip(\etastar)}_{P^{-1}}^2/\inner{\Fp(\theta)}{\theta-\etastar}$
($\theta\ne\etastar$). By \resref{rem:B1-effective-rate}, if I2 and I4 are replaced by the identities that define these
constants, the parameter part of the one-step bound~\eqref{eq:B1-effective-rate} contracts by a factor between
$1-\muEff^P/(2R_P^2)$ and $1-3\muEff^P/(16R_P^2)$, with $\muEff^P$ evaluated at $\theta_{t-1}$: the rate seen by the
proof is governed by $R_P^2/\muEff^P(\theta_{t-1})$.

\paragraph{Global constants of logistic regression.}
For $\ell_2$-regularized logistic regression, $\mu_P$ is explicit and $R_P^2$ is bounded on both sides
($\sigma_1\ge\dots\ge\sigma_d$ denote the eigenvalues of $\hat\Sigma$, not the per-sample curvatures
$\sigma_i'(\theta)$ of the lemma):
\restate{lem:B1-logistic}
Comparing the lower bound on $R_P^2$ with $\mu_P$ gives the central negative result:
\restate{prop:B1-nofreelunch}
The supremum in $R_P^2$ is attained on the decision boundary of a sample, where its data curvature is
$\frac14x_ix_i^\top$, whereas the infimum in $\mu_P$ is approached far from all of them, where $F''\to\lambda I$ is
isotropic; no single metric fits both, and with global constants the transferred Proposition~\ref{prop:59} never
certifies a faster rate than for $P=I$ on this problem. For the square loss, by contrast, $P=\hat\Sigma+\lambda I$ brings
$\kappaP$ into $[d_\infty(\lambda),d_\infty(\lambda)+1]$ (the ridge setting of~\cite{gonen2016solving};
Section~\ref{sec:B1-choice}). Here, with $X$ the matrix of rows $x_i^\top$,
$\mathrm{lev}_i(\lambda)=x_i^\top(X^\top X+n\lambda I)^{-1}x_i$ are the ridge leverage
scores~\cite{elalaoui2015fast,bach2013sharp}, $d_\infty(\lambda)=n\max_i\mathrm{lev}_i(\lambda)$ and
$d_{\mathrm{eff}}(\lambda)=\sum_i\mathrm{lev}_i(\lambda)=\tr(\hat\Sigma(\hat\Sigma+\lambda I)^{-1})$; by
\resref{lem:B1-leverage}, $d_{\mathrm{eff}}(\lambda)\le\min\{d_\infty(\lambda),d\}$ and $d_\infty(\lambda)<n$.

\paragraph{Trajectory-dependent constants.}
For logistic regression, any benefit of a metric must therefore come from the trajectory-dependent constants, which are
explicit in terms of the per-sample curvatures averaged on the segment $[\etastar,\theta]$:
\restate{prop:B1-traj-sc}
A sample with a large margin makes $m(\theta)$ small ($\sigma'(u)\approx e^{-\abs{u}}$) but has a small weight
in~\eqref{eq:B1-muEff-stat} unless $\theta-\etastar$ is aligned with it, so that $\muEff^P(\theta)$ is typically much
larger than $m(\theta)$; the global version of~\eqref{ineq:P-trajsc} is only $\mu_P=\lambda/(\sigma_1+\lambda)$, because
$\inf_\theta\min_i\sigma'(y_ix_i^\top\theta)=0$. The smoothness side is controlled by the leverage scores, by a
multiple of the effective dimension when they are incoherent:
\restate{ass:B1-coherence}
Since $\norm{x_i}_{(\hat\Sigma+\lambda I)^{-1}}^2=n\,\mathrm{lev}_i(\lambda)$ (\resref*{lem:B1-leverage}(i)),
\resref*{lem:B1-logistic} gives, for the statistical metric:
\restate{prop:B1-stat-global}
Dividing~\eqref{eq:B1-stat-RP} by~\eqref{ineq:P-trajsc} bounds the effective relative condition number:
\restate{cor:B1-kappa-traj}
In the last two bounds of~\eqref{eq:B1-kappa-traj-stat}, the eigenvalues of $\hat\Sigma$ appear only through
$d_\infty(\lambda)\le C\,d_{\mathrm{eff}}(\lambda)$, and $d_{\mathrm{eff}}(\lambda)$ counts the directions not killed by
the regularization, whatever the dynamic range of the spectrum. The price is the smallest segment curvature
$m(\theta)$, and on the instances of Section~\ref{sec:exp-B1} it dominates. There, as $\theta\to\etastar$,
$m(\theta)\to\min_iw_i$ with $w_i=\sigma'(y_ix_i^\top\etastar)$, and $\min_iw_i=7.4\cdot10^{-3}$ ($\alpha=1$) and
$6.1\cdot10^{-3}$ ($\alpha=2$), so that the first bound in~\eqref{eq:B1-kappa-traj-stat} tends to $1.7\cdot10^3$ and
$1.2\cdot10^3$: above $\kappa=10^3$, and above the upper bound in~\eqref{eq:B1-kappa-traj-I} for $P=I$ with the same
$m(\theta)$ ($968$ and $999$). On these instances this rigorous bound therefore predicts no gain of
$\hat\Sigma+\lambda I$.

\paragraph{Jacobi and Newton metrics; local constants.}
The Jacobi metric $P_{\mathrm J}=\diag(\hat\Sigma)+\lambda I$ costs $O(d)$ per iteration, like $P=I$, whereas a dense
$P$ costs $O(d^2)$ per iteration and a precomputation of order $nd^2+d^3$ (Section~\ref{sec:B1-cost}). With
$\hat\Sigma_{jj}>0$ the diagonal entries of $\hat\Sigma$, $x_{ij}$ the entries of $x_i$ and
$C_{\mathrm{corr}}=\diag(\hat\Sigma)^{-1/2}\hat\Sigma\diag(\hat\Sigma)^{-1/2}$ the correlation matrix of the features,
the constants of $P_{\mathrm J}$ compare with those of $\hat\Sigma+\lambda I$ as follows:
\restate{prop:B1-diag}
By~(iii), the constants of $P_{\mathrm J}$ are those of $\hat\Sigma+\lambda I$ up to the factor
$\operatorname{cond}(C_{\mathrm{corr}})$, in either direction: $P_{\mathrm J}$ removes the scale anisotropy of the
features but not their correlations, and nothing orders it with respect to $P=I$. The oracle Newton metric
$H=F''(\etastar)$, which requires $\etastar$, uses the curvatures $w_i\in(0,\frac14]$ at the optimum: with
$W=\diag(w_1,\dots,w_n)$ and $\hat\Sigma_w=\frac1nX^\top WX$ (eigenvalues $\sigma^w_1\ge\dots\ge\sigma^w_d$),
$H=\hat\Sigma_w+\lambda I$, and its constants involve the leverage scores
$\mathrm{lev}^w_i(\lambda)=w_ix_i^\top(X^\top WX+n\lambda I)^{-1}x_i$ of the weighted design and
$d^w_\infty(\lambda)=n\max_i\mathrm{lev}^w_i(\lambda)$:
\restate{prop:B1-newton}
The Newton metric makes strong convexity perfect near the optimum, but $R_H^2$ still pays the supremum over $\theta$:
sample $i$ enters it with its largest curvature $\frac14$ in the numerator and its curvature $w_i$ at the optimum in the
metric, hence the factor $1/w_i$ in~\eqref{eq:B1-newton-RP}. Near the optimum, the realized smoothness constant does
not pay this factor, and the realized strong convexity is a local curvature:
\restate{lem:B1-asymptotic}
Only these local constants account for the gains observed in Section~\ref{sec:exp-B1}. Whatever the direction of the
error, $\liminf_{\theta\to\etastar}\muEff^P(\theta)\ge\lmin(M_P)$ with $M_P=P^{-1/2}F''(\etastar)P^{-1/2}$, so that
$R_P^2/\lmin(M_P)$ bounds $\limsup_{\theta\to\etastar}\kappa^P_{\mathrm{eff}}(\theta)$. For $P=\hat\Sigma+\lambda I$,
$\lmin(M_P)$ is the smallest over directions $v$ of a weighted average of $w_1,\dots,w_n,1$ (weights $(x_i^\top v)^2/n$
and $\lambda\norm{v}^2$), not their minimum: on those instances it is $0.117$ ($\alpha=1$) and $0.109$ ($\alpha=2$),
$16$ and $18$ times $\min_iw_i$, and $R_P^2/\lmin(M_P)=166$ and $103$, against $R^2/\lmin(F''(\etastar))=576$ and $981$
for $P=I$ ($R_P^2$ being the upper bound in~\eqref{eq:B1-logistic-RP}, as in the experiments); over passes 15--20 the
measured $\kappa^P_{\mathrm{eff}}(\theta_{t-1})$ are about $84$ and $31$, against $460$ and $878$ for $P=I$ (geometric
means over states and seeds). These limits hold only as $\theta\to\etastar$, and their use as rates along a whole run
is heuristic; so used, the values $166$ and $103$ are below $\frac{9n}{16}=562.5$ ($n=1000$), and by
\resref{rem:B1-memory-bottleneck} the memory term $\frac1{3n}$ of the rate then binds. Through these limits,
Section~\ref{sec:B1-predictions} also predicts which step is the loosest late in a run (prediction (P2), heuristic
where it assumes the direction of the error): I2 for $P=I$, I4 with the global $\mu_P$ for the dense metrics
(Table~\ref{tab:B1-predictions-late} in Appendix~B1).

\paragraph{Known and new.}
Relative smoothness~\cite{lu2018relatively,bauschke2017descent}, Bregman-SAGA (whose rate has the same structure, with
constants $(\frac12,\frac18)$ instead of $(\frac13,\frac3{16})$~\cite[Corollary~1]{dragomir2021fast};
\resref*{rem:B1-dragomir}), the change of variables for a fixed metric~\cite{kolte2015accelerating,gonen2016solving}
and statistical preconditioning~\cite{hendrikx2020statistically} are known. To our knowledge,
\resref{prop:B1-nofreelunch} and \resref{cor:B1-kappa-traj} are not in the sources we reviewed in this form; the first
shows that, for logistic regression, the benefit of a fixed metric is entirely trajectory-dependent, and the second
does not certify the gains observed in Section~\ref{sec:exp-B1}, which only the local constants of
\resref{lem:B1-asymptotic} account for, and only heuristically.
